%% The manuscript body, shared by every binding of this paper.
%% The class, the front matter and the bibliography style belong to the
%% binding; everything a reader reads belongs here, once.
\section{Introduction}

A negative result usually survives inside a project as one sentence. Some
mechanism was tried, some number moved the wrong way, and the mechanism is
recorded as ruled out. The sentence outlives the arithmetic that produced it,
and by the time anyone reads it there is no way to tell whether the evidence
was strong, weak, or of a kind that could never have carried the sentence.

This paper takes one such sentence, recovers the arithmetic behind it, and
reports the registered paired run that was sized against that arithmetic. A
subject-driven personalisation study on a latent diffusion
backbone~\cite{rombach2022ldm,ruiz2023dreambooth} compared two ways of spending
the same low-rank adapter capacity~\cite{hu2022lora}: one adapter trained
across both attention types at once, against two adapters trained separately on
the cross-attention and self-attention projections and switched on together at
generation time~\cite{blora}. The composed arm came out with a higher
dispersion of identity fidelity across background prompts, \DisplayComposed{}
against \DisplayBase{}, and the project recorded the mechanism as killed.

The interesting part is not that the conclusion was too strong. It is why. The
runs themselves were paired in every sense that matters. All three arms on disk
were generated over the same \NPairs{} background prompts. The adapters were
trained at the same budget, the same learning rate, the same rank, the same
resolution and the same seed. The two composed adapters train \CrossParams{}
and \SelfParams{} parameters, which sum exactly to the \JointParams{} that the
joint arm trains as one, so the arms are matched on capacity rather than
merely similar in it. Every ingredient a paired test requires was present in
the experiment before anyone analysed it.

What removed the pairing was the endpoint. Identity fidelity was measured once
per background prompt, and then collapsed into its dispersion across prompts.
A dispersion across prompts is one number per arm. Two numbers admit no
pairing, no ranking by sign, and carry no correspondence for a test to
exploit. The comparison as registered therefore had no pairs at the moment it
was defined, whatever the runs had done, and the project's own record says so:
its pairing field records the comparison as unpaired, and the ladder entry that
cites the same two values downgrades them in prose. The caveat was written
down. The conclusion was never adjusted to match it.

The pattern is general. Endpoints of exactly this shape --- variance across
prompts, drift across scenes, spread across seeds, consistency across
conditions --- are ordinary in generative evaluation, and each of them
aggregates along the very axis a matched design controls. An endpoint that
aggregates along the matched axis converts a paired experiment into an unpaired
one with a single observation per arm, and additional generation at the prompt
level leaves that unchanged, because the count that matters has already been
reduced to one.

The paper contributes three things. First, it recovers the pairing from the
per-prompt records and reports what those \NPairs{} generations support when
analysed as the pairs they were run as. Second, it states a design rule that a
protocol can carry: the smallest probability an exact test can return at the
planned size is fixed by the number of pairs alone and is knowable before any
data exist. With \NPairs{} pairs that floor is \FloorAtFour{}, above
\Alpha{}, so no arrangement of the four historical results could have produced
a significant sign test; the historical comparison, taken on its own terms,
supports a ratio of dispersions anywhere from \DispLower{} to \DispUpper{}.
Third, it registers a paired protocol sized by that rule, with a pre-declared
kill rule, and reports its single execution: \RunPairs{} identity pairs across
\DesignConcepts{} concepts, declared before training and realised with
\RunFailures{} failures. The kill rule fired. The identity interval includes
zero and the registered test does not separate the arms. The protocol is in
the paper as the binding that was run, and the result is reported at the
strength that binding licenses.

\section{Related work}

This paper sits between two literatures that rarely meet: the engineering
literature on composing adapters, and the methodological literature on how
comparisons get registered and read. The case here is one where the second
determined what the first was allowed to conclude.

The technical setting is standard. Denoising diffusion began as a
nonequilibrium sampling construction~\cite{sohldickstein2015deep}, became a
practical generator once the reverse process was parameterised as a noise
prediction~\cite{ho2020ddpm,nichol2021improved,song2020ddim,song2021score,
dhariwal2021diffusion,karras2022edm,croitoru2023survey}, and its latent-space
successors made high-resolution synthesis
practical~\cite{rombach2022ldm,podell2023sdxl,goodfellow2014gan,
isola2017pix2pix}. Text-conditional systems then made the same process a
caption-to-image instrument~\cite{ramesh2021dalle,nichol2021glide,
ramesh2022dalle2,saharia2022imagen,schuhmann2022laion}; classifier-free
guidance, cross-attention editing and instruction-following edits made the
sampling process itself an object of
control~\cite{ho2022cfg,hertz2022prompt,zhang2023controlnet,meng2021sdedit,
brooks2023instruct}. Subject personalisation gave a way to bind a specific
subject to a model, first by inverting a word~\cite{gal2022textual} and then by
fine-tuning the weights~\cite{ruiz2023dreambooth}, and low-rank adaptation made
it cheap enough to train many such adaptations
separately~\cite{hu2022lora,han2023svdiff,kumari2023custom,ye2023ipadapter}.
Composing separately trained adapters rather than retraining jointly is an
active idea, including for implicitly decomposing an image into separable
components~\cite{blora}. The scores those papers print are usually a
distributional distance, a reference-free image--text similarity, or a learned
preference
proxy~\cite{heusel2017fid,radford2021clip,hessel2021clipscore,
kirstain2023pick,wu2023hps,xu2023imagereward}. Whether
composition costs anything relative to joint training is a reasonable question,
and the comparison we re-read was an attempt to answer it. The historical
reanalysis is an account of what that recorded attempt could and could not
support; the registered run then executed the protocol sized against that
limit.

The methodological literature is where the interesting part lies. That
flexibility in analysis can manufacture significance is well
established~\cite{simmons2011false,head2015phacking,ioannidis2005why,
openscience2015repro}, as is the subtler point that a single analysis chosen in
good faith, but contingent on the data seen, is enough to invalidate its own
probability~\cite{gelman2014garden}. Preregistration is the standard
remedy~\cite{nosek2018preregistration,munafo2017manifesto,pineau2021improving,
goodman2016mean}. The case here is a failure mode that remedy leaves open. The
comparison \emph{was} registered, and it was registered honestly, in advance.
What went wrong is that the registered endpoint --- a dispersion across
prompts, which is one number per arm --- discarded a pairing the runs had
actually produced. Preregistration protects against choosing the analysis
after seeing the data. It offers no protection against fixing, in advance, an
endpoint that destroys the structure of the experiment. We have not seen that
failure mode described as such, and it is ordinary: any summary that collapses
within-pair structure before the comparison does the same thing.

The inferential tools are old and well understood. The sign
test~\cite{dixon1946sign}, the signed-rank test~\cite{wilcoxon1945individual}
and the $t$ test~\cite{student1908probable} are the obvious candidates for the
recovered pairs, and the $t$ test's tolerance of non-normality at moderate
sample sizes is itself classical~\cite{box1953nonnormality}. The intervals we
report beside them are the ordinary bootstrap and binomial
constructions~\cite{efron1979bootstrap,clopper1934use}. What deserves more
attention than it gets is the discreteness of the exact tests at small $n$. With
\NPairs{} pairs the smallest two-sided probability the exact sign test can
return is \FloorAtFour{}, so no configuration of the data reaches \Alpha{}. This
is a stronger statement than low power. Power is the chance of detecting a real
effect of a stated size~\cite{cohen1992power}; an underpowered test may miss
one~\cite{button2013power}. This test had no way to report one. The
distinction matters because the remedy differs: power is improved by collecting
more data, whereas a floor above the threshold means the test was the wrong
instrument regardless of what the data showed. Calls to stop dichotomising at a
fixed threshold~\cite{amrhein2019retire,wasserstein2016asa} are usually made
about results that cleared it or narrowly missed; the same argument applies with
more force to a test whose attainable minimum sits above the line before any
data arrive.

\section{Methods}
\subsection{What was already on disk}

Nothing in this study was generated for it. Every quantity comes from records
written by runs that finished before this analysis began, and each record is
bound to the manuscript by its SHA-256 digest; the figure code re-hashes each
file before reading it and stops the build on any mismatch. No image was
generated, no adapter was trained, and no evaluation was launched.

The records describe a subject-driven personalisation study on a latent
diffusion backbone. A concept is taught to the backbone through low-rank
adapters on the attention projections, and the resulting model is asked to
place that concept in a set of background scenes. Three arms exist on disk. The
\emph{jointly trained} arm puts one adapter on both attention types at once.
The \emph{composed} arm trains two adapters separately --- one on the
cross-attention projections, one on the self-attention projections --- and
switches both on at generation time. A third arm switches the self-attention
adapter off and leaves the cross-attention adapter alone.

\subsection{The comparison as it was recorded}

The registered comparison is between the jointly trained arm and the composed
arm on a single endpoint: the dispersion of identity fidelity across background
prompts. Each arm generated one image per prompt for \NPairs{} prompts, and the
dispersion is the standard deviation of the \NPairs{} resulting fidelities.
It is a population standard deviation, dividing by the number of prompts rather
than by one fewer; the figure code recomputes it from the per-prompt values and
refuses to proceed unless it reproduces the recorded number exactly, because
the entire comparison is between two such numbers.

The project recorded the comparison's evaluation pairing as
\texttt{\PairingLabel{}} and recorded that the pre-registered criterion
\KillFires{} on it. It also recorded, in the same file and in the ladder entry
that cites it, that this is an unpaired historical result rather than a paired
measurement. The present paper takes that caveat as its starting point and asks
what the same bytes do support.

\subsection{What the arms share}

Three configuration records, one per adapter, are bound here. The figure code
checks that the three agree on the backbone checkpoint, the training budget,
the learning rate, the adapter rank, the resolution, the training seed and the
instance prompt, and stops the build if any of them differ. It also checks
that the trainable parameter counts of the two composed adapters sum exactly to
the count the joint arm trains, which is the sense in which the two arms are
matched on capacity rather than merely similar.

\subsection{Reconstructing the pairs}

The per-prompt values underlying each dispersion are recorded, one row per
background prompt, with the prompt text as the row key. Rows are joined across
arms on that text rather than on position, so a reordering cannot silently pair
the wrong generations. Two endpoints are available per row: identity fidelity,
which the registered comparison aggregated, and prompt fidelity, which it did
not use.

A pair is one background prompt, and the difference is the composed arm's value
minus the jointly trained arm's value at that prompt. This recovers the
correspondence the dispersion discarded; it does not recover a pairing for the
registered endpoint itself, which cannot have one, because a dispersion across
prompts is a single number per arm.

\subsection{Tests, and what each of them can reach}

Three tests are reported on the same \NPairs{} pairs. The exact two-sided sign
test counts how many differences fall each way; ties would be dropped by
convention and are refused here instead, because with four pairs one discarded
tie changes the answer. The Wilcoxon signed-rank test on the same pairs is
reported beside it. The paired $t$ test is reported third, with its interval,
and is the only one of the three that assumes a distributional shape.

For the recorded comparison, which has no pairs, we report the interval its own
design supports: the ratio of the two dispersions with the confidence interval
an $F$ distribution gives at \NPairs{} generations per arm. That interval
assumes the two arms are independent, which they are not --- they share the
prompts. The assumption is stated rather than corrected because it is the
assumption the unpaired comparison implicitly makes, and the interval it yields
is already wide enough to settle the question this paper asks of it.

\subsection{The attainable probability of a test, before any data}

The exact sign test at $n$ pairs can return only the probabilities its
binomial reference distribution contains. The smallest two-sided value is
reached when every pair falls the same way and equals $2 \cdot 2^{-n}$; if $k$
pairs are allowed to fall the other way the smallest attainable value is
$2 \sum_{i \le k} \binom{n}{i} 2^{-n}$. These quantities depend on $n$ alone
and can be read before the experiment exists. They are what the pre-registered
design is sized against.

\subsection{How the protocol's size was chosen}

The protocol below is derived rather than picked. The background prompts are
the recorded grid and are not changed. The source design requires a second
concept, so that a conclusion about the mechanism can be distinguished from a
conclusion about one subject, and it requires the conclusion to be readable per
concept; each concept therefore needs enough pairs on its own. The number of
generation seeds is the smallest that leaves each concept able to reach
\Alpha{} with at least one pair falling the wrong way, which follows directly
from the attainable-probability expression above. Everything else in the
protocol table is a consequence of that choice.

\subsection{What is deliberately not computed}

The reference comparison this line of work would eventually be measured
against was never produced, and its tier is recorded as \BlockedTier{}. No
value is imputed for it, no substitute metric is introduced in its place, and
the two perceptual similarity measures the wider literature reports are
recorded on disk as not computed and are not computed here either. The recorded
gate list for the unrun lane stands at \AuthorisationState{}, and the paired
evaluation this paper specifies is not executed: the deliverable is the
specification.


\section{The pairing was in the runs}

Fig.~\ref{fig:capacity} and Table~\ref{tab:arms} describe what was trained.
The arm trained as a single adapter puts low-rank updates on both attention
types at once and trains \JointParams{} parameters. The composed arm trains
two adapters separately, \CrossParams{} parameters on the cross-attention
projections and \SelfParams{} on the self-attention projections, and switches
both on at generation time. Those two counts sum to \JointParams{} exactly, and
the figure code stops the build if they ever differ, because the paper's
framing depends on it: the comparison is between two ways of spending one
parameter budget, not between a larger model and a smaller one.

The three training records also agree on the backbone checkpoint, a budget of
\TrainSteps{} steps, learning rate \LearningRate{}, adapter rank \LoraRank{},
resolution \Resolution{}, seed \TrainSeed{} and the instance prompt. Agreement
on all seven is checked rather than assumed. A comparison in which the arms
differ on any of them is a different paper, and it is easier to verify the
seven than to remember them.

Generation was matched too. Each arm produced one image per background prompt
over the same \NPairs{} prompts, and the per-prompt values survive on disk with
the prompt text as the private row key. Rows are joined across arms on that key
rather than on position, so a reordering of the records leaves the pairing
intact; the public table and figures use stable P1--P4 row codes.

Fig.~\ref{fig:collapse} shows what happened next. Each line in the figure is
one background prompt followed across the three arms, and the number printed
beneath each arm is what the recorded comparison actually used: the standard
deviation of that arm's \NPairs{} identity fidelities. The lines are the
experiment. The three numbers underneath are the analysis. Going from the first
to the second is where the pairing was lost, and it is a one-way step: nothing
downstream of those three numbers can recover which prompt produced which
value.

\begin{figure}[tbp]
\centering
\includegraphics[width=\linewidth]{fig1_capacity.pdf}
\caption{\CapCapacity}
\label{fig:capacity}
\end{figure}

\begin{table}[tbp]
\centering
\caption{\CapArms}
\label{tab:arms}
\small
% Generated by the figure script from hash-bound evidence. Do not hand-edit.
\begin{tabular}{lrrr}
\toprule
Adapter & Trainable parameters & Wall clock (s) & Final loss \\
\midrule
Jointly trained & 3,188,736 & 119.9 & 0.0301 \\
Cross-attention adapter & 1,591,296 & 121.8 & 0.0359 \\
Self-attention adapter & 1,597,440 & 128.1 & 0.0311 \\
\midrule
Two adapters combined & 3,188,736 & 249.9 & --- \\
\bottomrule
\end{tabular}

\end{table}

\begin{figure}[tbp]
\centering
\includegraphics[width=\linewidth]{fig2_collapse.pdf}
\caption{\CapCollapse}
\label{fig:collapse}
\end{figure}

\section{What the recorded comparison supports on its own terms}

Before recovering anything, it is worth asking what the registered comparison
would support if one simply accepted it. Two dispersions, each estimated from
\NPairs{} generations, differ by a factor of \SDRatioComposed{}. Treating the
arms as independent samples --- which is what an unpaired comparison of two
dispersions assumes --- the ratio of variances follows an $F$ distribution, and
the interval that yields is shown in Fig.~\ref{fig:dispersion}.

The interval runs from \DispLower{} to \DispUpper{}, at $F=\DispF{}$ and a
two-sided probability of \DispP{}. It contains no difference. It contains a
halving. It contains a near fivefold increase. The observed direction is inside
it and so is the opposite direction, and the width is not a consequence of
noisy measurement: it is what an estimate of dispersion from \NPairs{} values
is worth. Dispersion estimates are among the least stable quantities one can
compute from a small sample, and they are additionally sensitive to
departures from normality in a way that comparisons of means are
not~\cite{box1953nonnormality}.

The independence assumption is worth naming rather than fixing. The two arms
share the background prompts, so they are not independent, and the interval
above is therefore not the right interval for these data. We report it because
it is the interval the registered comparison implicitly claims, and because it
is already wide enough to answer the question we are asking of it. An interval
that still contains a halving settles the direction in neither sense.

The third arm makes the same point from the other side. Switching the
self-attention adapter off leaves a dispersion of \SDIdentityOnly{}, a ratio of
\SDRatioIdentityOnly{} to the jointly trained arm, with an interval from
\DispIdonLower{} to \DispIdonUpper{} and a two-sided probability of
\DispIdonP{}. Read as the registered comparison reads its own endpoint --- less
dispersion is better --- the best arm on disk is the one that discards half the
trained capacity, and it beats the joint arm by a wider margin than the
composed arm loses to it. We read that as a property of the endpoint rather
than as a finding about adapters: a dispersion endpoint rewards an arm for
varying less, including for varying less because it has less to vary with,
and that is why a dispersion endpoint needs a level endpoint printed beside
it. On level, as shown below, the same arm is indistinguishable from the joint
arm.

\begin{figure}[tbp]
\centering
\includegraphics[width=0.9\linewidth]{fig3_dispersion.pdf}
\caption{\CapDispersion}
\label{fig:dispersion}
\end{figure}

\section{The same generations, analysed as pairs}

The per-prompt values are on disk, so the correspondence the dispersion
discarded can be put back. Table~\ref{tab:pairs} gives every value and every
difference; Fig.~\ref{fig:paired} plots the differences.

On identity fidelity, the endpoint the registered comparison aggregated, the
composed arm is lower at \IdentityDown{} of \NPairs{} prompts. The exact
two-sided sign test~\cite{dixon1946sign} returns \IdentitySignP{}, the Wilcoxon
signed-rank test~\cite{wilcoxon1945individual} on the same pairs returns
\IdentityWilcoxP{}, and the mean paired difference is \IdentityDiff{} with a
95\% interval from \IdentityLower{} to \IdentityUpper{}. The interval contains
zero and both of its ends are larger in magnitude than the difference itself.

On prompt fidelity, which is recorded in the same files and was outside the
registered endpoint, the composed arm is lower at all \PromptDown{} prompts.
The sign test returns \PromptSignP{}, Wilcoxon returns \PromptWilcoxP{}, and
the mean paired difference is \PromptDiff{} with an interval from
\PromptLower{} to \PromptUpper{} that excludes zero. The paired $t$
test~\cite{student1908probable} on that endpoint returns \PromptTP{}.

That last number is the only quantity in this historical half below \Alpha{},
and its place in the paper is fixed by the protocol rather than by its size.
It was not the registered endpoint. It was found by looking at a second column
that happened to be in the same files, and reading it as the result would be
the exact manoeuvre that makes small studies unreliable: an endpoint chosen
after the analysis, reported at the threshold it happens to
clear~\cite{simmons2011false,gelman2014garden}. It also rests on a normality
assumption that four points can neither support nor
contradict~\cite{amrhein2019retire}. What the direction deserves is a line in a
protocol, and it has one: prompt fidelity was written into the design below as
a pre-registered secondary endpoint, and it was then tested on pairs that had
no part in selecting it. It is reported from the executed run at that
secondary rank.

The third arm is consistent with all of this and adds one detail. Switching the
self-attention adapter off moves identity fidelity down at \IdonDown{} of
\NPairs{} prompts, a sign-test probability of \IdonSignP{}, which is to say the
arm with the lowest dispersion on disk is indistinguishable from the joint arm
on the level of the same quantity. The two readings of that arm --- best by
dispersion, tied on level --- come from the same eight numbers.

\begin{table}[tbp]
\centering
\caption{\CapPairs}
\label{tab:pairs}
\small
% Generated by the figure script from hash-bound evidence. Do not hand-edit.
\begin{tabular}{lrrrrrr}
\toprule
& \multicolumn{3}{c}{Identity fidelity} & \multicolumn{3}{c}{Prompt fidelity} \\
\cmidrule(lr){2-4}\cmidrule(lr){5-7}
Background prompt & Joint & Composed & Difference & Joint & Composed & Difference \\
\midrule
P1 & 0.8328 & 0.8126 & -0.0202 & 0.2715 & 0.1678 & -0.1036 \\
P2 & 0.8762 & 0.7952 & -0.0810 & 0.2372 & 0.1735 & -0.0637 \\
P3 & 0.8260 & 0.8625 & +0.0365 & 0.2173 & 0.1826 & -0.0347 \\
P4 & 0.8579 & 0.8348 & -0.0231 & 0.2710 & 0.2096 & -0.0614 \\
\midrule
Dispersion across prompts & 0.0200 & 0.0252 & --- & --- & --- & --- \\
\bottomrule
\end{tabular}

\end{table}

\begin{figure}[tbp]
\centering
\includegraphics[width=\linewidth]{fig4_paired.pdf}
\caption{\CapPaired}
\label{fig:paired}
\end{figure}

The two per-prompt coordinates can be kept separate without choosing a new
endpoint. Fig.~\ref{fig:endpoint-map} places the composed and cross-attention-only
changes relative to the jointly trained arm in that two-dimensional space. All
\PromptDown{} composed-arm prompt-fidelity changes point in the same direction,
whereas \IdentityDown{} identity changes are lower and \IdentityHigher{} are
higher (with \IdentityTies{} exact ties). The map is a descriptive diagnostic
of the recorded rows: it keeps the two coordinates separate and adds no pooled
trade-off, correlation or causal path. The prompt-fidelity coordinate is
recorded in the same files and sat outside the historical registered endpoint.

\begin{figure}[tbp]
\centering
\includegraphics[width=0.66\linewidth]{fig6_tradeoff.pdf}
\caption{\CapEndpointMap}
\label{fig:endpoint-map}
\end{figure}

\section{What the historical test could reach}

Neither sign test above reaches \Alpha{}, and the usual reading of that is
that the effect is absent or the data are noisy. Here the explanation is
structural.

An exact sign test at $n$ pairs draws its probabilities from a binomial
reference distribution, and that distribution is fixed by $n$ before any datum
exists. The smallest two-sided value it contains is reached when every pair
falls the same way and equals $2 \cdot 2^{-n}$. With \NPairs{} pairs that value
is \FloorAtFour{}, which is above \Alpha{}. Every possible outcome of the
registered experiment, including perfect unanimity in either direction, returns
a probability at or above \FloorAtFour{}. The test had no significant outcome
available to it.

Fig.~\ref{fig:floor} plots this floor against the number of pairs, for
unanimity and for designs that allow one or two pairs to fall the wrong way.
Reaching \Alpha{} at all requires \MinNAny{} pairs, and then only on unanimity:
a single discordant pair at that size puts the result back above the threshold.
Tolerating one discordant pair requires \MinNOne{}, and tolerating two requires
\MinNTwo{}.

Three properties of these numbers are worth stating plainly, because together
they make the quantity cheap to use. They depend on the number of pairs alone,
so no assumption about effect size, variance or distributional shape enters.
They are computable before the experiment exists, from the protocol rather than
from data. And they are a hard ceiling on what the analysis can conclude, not a
soft expectation like power: a power calculation says an effect may be missed,
whereas this says which answers a specific test statistic has available to it.

The practical consequence is a line that a protocol can carry and a reader can
check: state, before running, the smallest probability the primary test can
return at the planned sample size. If that number is above the threshold the
study intends to use, the study has already decided its own outcome, and the
choice is to enlarge it, to change the test, or to declare in advance that it
is descriptive. Underpowered studies are a long-recognised problem in
experimental fields~\cite{button2013power}; the case here is the sharper corner
of it, where the reachable set of answers excludes the one the study is
looking for.

\begin{figure}[tbp]
\centering
\includegraphics[width=\linewidth]{fig5_floor.pdf}
\caption{\CapFloor}
\label{fig:floor}
\end{figure}

\section{The paired experiment, locked and then executed}

Table~\ref{tab:protocol} is the protocol that was locked before the run, and
Fig.~\ref{fig:protocol-flow} draws the same specification as a sequence, with
the declaration time and the realised pair count read from the bound
pre-declaration. Every number in the table is derived from the arithmetic
above rather than chosen, and the derivation is worth stating so that a
reader can disagree with a step instead of with the total.

The pairing unit is one concept, one background prompt and one generation seed,
with the two arms generated under that triple and differing only in whether the
adapters were trained jointly or separately. Naming the seed inside the unit is
what makes the pair a pair; the historical comparison is recorded as unpaired
precisely because no such unit was ever written down.

The background prompts are the recorded grid, unchanged at \NPairs{}. Keeping
them fixed costs nothing and means the executed run can be read against the
historical half.

Concepts are \DesignConcepts{}. Only \NConceptsRun{} concept existed on disk
when the protocol was locked, so nothing then distinguished a conclusion about
the mechanism from a conclusion about one subject, and the source design had
already required a second concept for that reason. The registered run
realises both pre-specified concepts. The primary inferential decision is
pooled across the two pre-specified concepts at \Alpha{}; concept-specific
estimates are shown as descriptive strata and sit outside the confirmatory
test family.

Seeds per prompt are \DesignSeeds{}, and this is the only number with any
freedom in it. The pooled design needs at least \MinNOne{} pairs to reach
\Alpha{} while tolerating a single discordant pair, and \DesignSeeds{} seeds
across \NPairs{} prompts is the smallest multiple of the prompt grid that clears
it. That gives \DesignPerConcept{} pairs per concept and \DesignTotal{} in
total, tolerating \DesignTolerance{} discordant pairs within a concept and
\DesignPooledTolerance{} pooled. Requiring unanimity would have been cheaper
and is the wrong trade: a design that one contrary pair overturns is a design
that will be re-run until none appears.

The primary endpoint is the per-pair difference in identity fidelity, tested by
one pooled exact two-sided sign test at \Alpha{} across the two concepts, with
concept-specific estimates reported descriptively. Wilcoxon on the same pairs
and prompt fidelity are secondary checks, kept outside the confirmatory
family. Effect sizes and intervals are reported for all of them regardless of
outcome.

The training budget is \TrainSteps{} steps, which is what the arms being
re-tested were trained at. The source design named a shorter budget of
\ProxySteps{} steps, borrowed from a separate proxy arm, and we depart from it
deliberately. A protocol that re-tests a comparison must train at the budget of
the arms it re-tests; at any other budget it answers a different question and
leaves the original untouched. The proxy arm gives a second reason to keep
away from it: two records of that same short run, agreeing on entry point,
parameter count and final loss, report its wall clock as \ProxyWallRun{} and
\ProxyWallCloseout{} seconds respectively. The discrepancy is a few seconds
and of no consequence in itself, and it is exactly the kind of thing one
should notice before sizing a budget from a record.

The cost follows. Training is \DesignTrainRuns{} runs, three adapters for each
of \DesignConcepts{} concepts. The three arms on disk took \TrainSecondsOne{}
seconds of wall clock between them, so the training for both concepts is
about \TrainSecondsAll{} seconds at the recorded budget. Generation is
\DesignGenerations{} images, two arms across \DesignTotal{} pairs. This is a
small experiment, and that is the point: the original question was left open
by its endpoint, not by the cost of the runs.

\begin{table}[tbp]
\centering
\caption{\CapProtocol}
\label{tab:protocol}
\small
% Generated by the figure script from hash-bound evidence. Do not hand-edit.
\begin{tabular}{ll}
\toprule
Element & Locked declaration (executed once) \\
\midrule
Pairing unit & one concept, one background prompt, one generation seed \\
Arms within a pair & jointly trained adapter versus the two composed adapters \\
Concepts & 2 \\
Background prompts & 4, the recorded grid, unchanged \\
Generation seeds per prompt & 3 \\
Pairs per concept & 12 \\
Pairs in total & 24 \\
Training budget & 500 steps, the budget of the arms being re-tested \\
Primary endpoint & per-pair difference in identity fidelity \\
Primary test & two-sided exact sign test at $\alpha=0.05$ \\
Secondary endpoint & per-pair difference in prompt fidelity \\
Secondary test & Wilcoxon signed rank on the same pairs \\
Discordant pairs tolerated per concept & 2 \\
Discordant pairs tolerated pooled & 6 \\
Training runs required & 6 \\
Generations required & 48 \\
\bottomrule
\end{tabular}

\end{table}

\begin{figure}[tbp]
\centering
\includegraphics[width=\linewidth]{fig7_protocol.pdf}
\caption{\CapProtocolFlow}
\label{fig:protocol-flow}
\end{figure}

\section{The registered run}

The protocol above was declared before any training
(digest \PredeclarationShaShort{}) and then executed once. All
\RunPairs{} pairs were realised, the failure ledger has
\RunFailures{} entries, and Fig.~\ref{fig:paired-run} shows every paired
difference.

The registered primary endpoint is the per-pair difference in identity
fidelity, tested by one pooled exact two-sided sign test at \Alpha{}. The
composed arm was lower on \RunIdentityLower{} of \RunPairs{} pairs and
higher on \RunIdentityHigher{}. The sign-test probability is
\RunIdentitySignP{}. The mean paired-difference interval runs from
\RunIdentityCiLo{} to \RunIdentityCiHi{} and includes zero. The confirmatory
test does not separate the arms. At \RunPairs{} pairs the registered test
rejects only when at least \SignCriticalCount{} pairs fall the same way, so it
had \SignPowerTarget{}\% power only for a true composed-lower proportion of at
least \SignPowerEighty{}; the observed proportion \SignObservedProp{} (exact
95\% interval \SignObservedPropLo{} to \SignObservedPropHi{}) lies inside the
region the design left unresolved, so the result is a non-separation at this
size rather than evidence of equality.

Prompt fidelity is the registered secondary endpoint. The composed arm was
lower on \RunPromptLower{} of \RunPairs{} pairs at a sign-test probability
of \RunPromptSignP{}. The protocol fixed that rank before the run and the
paper keeps it: the count is reported as the declared secondary, it is not
promoted to a confirmatory claim, and the identity conclusion rests on the
primary endpoint alone.

Concept strata are descriptive. On identity fidelity the composed arm was
lower on \RunDogIdentityLower{} of the twelve dog pairs and higher on
\RunClockIdentityHigher{} of the twelve clock pairs. The two strata are read
as description, outside the confirmatory family.

The historical dispersion ratio on the \NPairs{}-prompt grid was
\SDRatioComposed{}. On the executed run the recomputed composed-over-joint
dispersion ratios are \RunDispersionDog{} (dog) and \RunDispersionClock{}
(clock). Both are below one, the reverse of the historical point, and both
are reported as description.

The pre-registered kills settle as follows. Kill 1 fires: the identity
interval contains zero, so the claim is that no separation was observed on
the registered primary. Kill 2 does not fire: the pooled identity count is
not a significant reversal of the historical direction. Kill 3 is
descriptive only: the two concept strata differ in sign, and that
difference is read as description rather than as a confirmatory close.

\begin{figure}[tbp]
\centering
\includegraphics[width=\linewidth]{fig8_paired_run.pdf}
\caption{\CapPairedRun}
\label{fig:paired-run}
\end{figure}

The paired run can also be read without collapsing either endpoint. Fig.~\ref{fig:concept-slopes}
joins the two arms within each concept, prompt and generation seed. Identity
is the primary endpoint and prompt fidelity remains secondary; the slopes are
descriptive views of the same \RunPairs{} rows and add no test.

\begin{figure}[tbp]
\centering
\includegraphics[width=0.8\linewidth]{fig9_concept_slopes.pdf}
\caption{\CapConceptSlopes}
\label{fig:concept-slopes}
\end{figure}

\section{Discussion}

The specific findings are three. One recorded mechanism result rested on a
comparison whose design admitted no paired test. The strongest paired analysis
of those same four generations leaves the question open. And the registered
protocol sized against that floor, once executed, reports no identity
separation at \RunPairs{} pairs. Five points travel beyond this case, because
none of them is specific to adapters or to diffusion.

The first is that an endpoint can undo a design. The experiment was matched on
background prompt. The endpoint aggregated across background prompt. Those two
sentences are individually unremarkable and jointly fatal, and nothing in the
usual checklist catches the combination: the arms were controlled, the seeds
were fixed, the capacity was matched, and the analysis still had one number per
arm. The check that would have caught it is one question asked at design time
--- does the primary endpoint preserve the unit the design matched on? --- and
if the answer is no, either the endpoint or the claim has to change. Endpoints
that aggregate across a matched axis are common in generative evaluation, where
stability, drift and consistency are all naturally written as spreads.

The second is to print what the test can reach before running it. The
attainable floor of an exact test is a property of the design, costs one line
to compute, and would have flagged the original comparison as unable to reach
its own threshold before a single image was generated. Pre-registration
already asks for the endpoint, the test and the sample
size~\cite{nosek2018preregistration}; adding the smallest probability the test
can return at that size makes the combination self-checking, and it is a
stronger statement than a power calculation because it depends on no guessed
effect.

The third is that recording a caveat is different from propagating it. The
pairing status was recorded, in a machine-readable field, in the same file as
the conclusion, and the ladder entry that cites the file downgrades the result
in prose. The qualification was never lost. It simply did not reach the
sentence that travelled. This is the failure a schema can actually prevent: a
conclusion field whose strength is checked against a pairing field in the same
record is a cheap constraint, and the alternative is to rely on every future
reader noticing a caveat two files away.

The fourth is that a dispersion endpoint needs a level endpoint beside it.
Read by dispersion alone, the best arm on disk is the ablation that switches
half the composed capacity off. Read by level, that arm is indistinguishable
from the joint arm. Both readings come from the same values, and reporting
only the first would have suggested an intervention --- drop the
self-attention adapter --- that the second does not support. Any evaluation
whose headline is a spread is recommending the arm that moves least, and
should say what happens to the quantity that is supposed to be high.

The fifth is that re-analysis is cheap and bounded. The historical half of
this paper is arithmetic over records that already existed; it trained nothing
and generated nothing. What it can deliver is a corrected strength for a
claim: the strongest paired analysis of those four pairs leaves the question
open, which is why the registered protocol was executed rather than left as a
specification. The evidence for a claim has to come from a run. The registered
run supplied that evidence, and at this size it reports no identity
separation.

\section{Scope of the claims}

The design licenses one comparative claim, and this paper makes exactly that
one: on the registered primary endpoint, a pre-declared paired sign test at
\RunPairs{} pairs does not separate the composed arm from the jointly trained
arm on identity fidelity. Whether composing two separately trained adapters
is better or worse than training the same parameters jointly is left open by
that result, in both directions.

Prompt fidelity holds the rank the protocol gave it. On the executed run it is
the declared secondary endpoint and is reported as such. The historical
prompt-fidelity $t$ test that cleared \Alpha{} remains an unregistered analysis
of four pairs, reported for completeness and set aside as a basis for any
claim.

The endpoints used are the two recorded on disk, identity fidelity and prompt
fidelity. The record states that the two perceptual similarity measures the
wider literature reports were not computed, and this paper works with the
recorded endpoints only. The reference procedure at full length was never
produced, so its cell stays empty: the paper reports its recorded status,
imputes no value for it, introduces no substitute, and presents nothing here
as an approximation of it.

The concept strata are descriptive. The registered run includes both declared
concepts; the dog and clock identity counts differ in sign and are read as
description rather than as a second test family.

The result reported is the single execution authorised under the bound
pre-declaration, digest \PredeclarationShaShort{}. The protocol admits no
re-run and no unregistered variant, and none is reported.

The short-budget proxy arm enters only as a record-consistency observation.
It was trained at \ProxySteps{} steps rather than \TrainSteps{}, it is a
different arm from the three compared here, and its only role in this paper is
the note that two records of it disagree about its wall clock.

The method that originated the attention-type split is cited for the split
alone. This paper re-analyses one small local comparison whose design borrows
that split; it does not evaluate the method, does not reproduce it, and
reports no number about it.

\section{Conclusion}

A recorded mechanism result in a diffusion personalisation project was traced
back to the artifacts it was computed from. The runs behind it were paired:
three arms over the same \NPairs{} background prompts, adapters trained at the
same budget, rank, resolution, learning rate and seed, and a composed arm whose
two adapters train exactly the \JointParams{} parameters the joint arm trains.
The pairing was discarded by the endpoint rather than by the experiment,
because a dispersion across prompts is one number per arm.

Taken as recorded, the comparison supports a ratio of dispersions from
\DispLower{} to \DispUpper{}, which excludes nothing of interest. Taken as
pairs, the composed arm is lower on identity fidelity at \IdentityDown{} of
\NPairs{} prompts at a sign-test probability of \IdentitySignP{}, and lower on
prompt fidelity at all \PromptDown{}, at \PromptSignP{}. Neither reaches
\Alpha{}, and the reason is structural: at \NPairs{} pairs the smallest
attainable two-sided probability of an exact sign test is \FloorAtFour{}. The
one analysis that does cross \Alpha{} uses an endpoint outside the registration
and a distributional assumption that four points leave untested; the protocol
gave that endpoint a secondary place and tested it on pairs that had no part
in selecting it.

The historical finding restated at the strength its evidence carries is that
no separation was observed in one small comparison whose registered endpoint
admits no paired test. The experiment that could settle that comparison was
then executed as declared: \RunPairs{} identity pairs across
\DesignConcepts{} concepts, pairs defined at generation time by concept,
prompt and seed. The confirmatory sign test returned
\RunIdentitySignP{}, and the identity interval includes zero. The protocol
carries, in its sizing derivation rather than as a row of
Table~\ref{tab:protocol}, the line whose absence is the subject of the
historical half of this paper: what its test can return at the chosen size
was fixed before any of it happened, by the attainable-probability rule of the
Methods and the sizing read from Fig.~\ref{fig:floor}, and the table lists the
design cells sized against that rule. The historical grid's own floor,
\FloorAtFour{} at \NPairs{} pairs, is the case that made the line necessary;
at the \RunPairs{} pairs realised, the same rule fixes the smallest attainable
two-sided probability at \SignFloorRealised{}, below \Alpha{}, and that value
is printed here so the claim does not rest on the reader's arithmetic.
The scoped claim is that this pre-declared test does not separate the arms on
identity fidelity.

\section*{Data and code availability}

The recorded comparison, the per-prompt fidelity tables underlying it, the
three training records, the recorded status of the reference comparison and of
the authorisation gate, and the figure code that produces every number in this
manuscript accompany the manuscript source. Each artifact is recorded under an
evidence ID in a manifest, and the figure code verifies the bound bytes before
reading them, so the build stops rather than reporting stale numbers. The
registered run generated \DesignGenerations{} images locally; those images,
and the adapter weights, are retained under a gitignored outputs directory and
are not redistributed. The Stage-1 archive cited next is the earlier
protocol-stage deposit, made before the run, and is distinct from this
executed manuscript. The backbone checkpoint is a third-party release
identified in the manifest by name rather than copied, and the subject images
are from a publicly distributed subject-driven generation set held under its
own licence.
%% Generated by scripts/release_targets.py from the release ledger.
%% Identifiers below are ledger-checked. The sentence is a Task-5 honesty
%% pass: 22647020 is the Stage-1 unrun-protocol archive, not this executed draft.
%% Ledger state: github=CREATED, zenodo=DEPOSITED
\ifdefined\ANONYMOUS The Stage-1 protocol archive has a persistent identifier; that deposit is the previous unrun-protocol paper, not this executed draft, and the addresses are withheld from this review copy.\else The Stage-1 protocol archive is under the persistent identifier \href{https://doi.org/10.5281/zenodo.22647020}{10.5281/zenodo.22647020}; that deposit is the previous unrun-protocol paper, not this executed draft. A development repository is listed at \url{https://github.com/PeterPonyu/discarded-pairing-reanalysis}.\fi


\section*{Ethics approval and consent to participate}

Not applicable: the study involves no human participants and no animal
subjects, and no personal data are processed. The historical half is a
secondary analysis of numeric records produced by earlier computational runs.
The registered half generated \DesignGenerations{} images locally; they are
retained under a gitignored outputs directory, are not distributed, and none
is reproduced in the manuscript. The subjects depicted in the personalisation
runs are from a publicly distributed set assembled for subject-driven
generation research. Generative image models raise well-documented concerns
around consent for training imagery and around synthesis of identifiable
subjects; this paper operates on that same publicly distributed subject set
under the same terms.

\section*{Consent for publication}

Not applicable. No identifiable individual's data or likeness is included,
and no generated image is reproduced in the manuscript.

%% Funding and competing-interest wording adopted from the owner's packaging
%% script (scripts/build_springer_submission_package.py), shipped in every
%% packet since 2026-08-28 (ruling R16, 2026-09-09). FUNDING_AND_CONFLICTS
%% stays OPEN for the owner's final read.
\section*{Competing interests}

The author declares that they have no competing financial or non-financial
interests that are directly or indirectly related to the work submitted for
publication.

\section*{Funding}

No external funding was received for this study.

\section*{Author contributions}

\ifdefined\ANONYMOUS
The sole author performed conceptualization, methodology, software, validation,
formal analysis, investigation, visualization, writing---original draft, and
writing---review and editing.
\else
Zeyu Fu: conceptualization, methodology, software, validation, formal analysis,
investigation, visualization, writing---original draft, and writing---review
and editing.
\fi

%% TODO-OWNER: Acknowledgements. No factual content is available to the
%% editorial pass (no funder, no named contributor, no shared facility recorded
%% in the ledger). Machine Learning (Springer) treats the section as optional.
%% If the owner wants one, uncomment and fill:
%% \section*{Acknowledgements}
%% TODO-OWNER

\section*{Declaration of generative AI use}

Generative AI tools assisted with code preparation and language editing. Every
number printed above is emitted programmatically from artifacts tied to the
manuscript by cryptographic evidence binding rather than transcribed into
prose. Structural claims --- including the adapter partition, shared training
settings, and recomputed per-prompt dispersions --- are enforced by build-time
checks that stop the build. Responsibility for the content rests with the
author.
